\section{The ambient projection onto a finite span}

\begin{theorem}[Projection matrices and actual ranges]
    \label{thm:projection_matrix_range}
    \lean{Matrix.range_toEuclideanLin_eq_map_range_mulVecLin,
        Matrix.projectionMatrix_range_of_isStarProjection}
    \leanok
    Let $A:\C^n\to\C^m$ be a rectangular matrix. The standard coordinate
    identification carries its ordinary column range exactly to the range
    of its Euclidean-space action. If $P$ is an orthogonal projection
    matrix, then the matrix of the orthogonal projection onto its actual
    range is $P$ itself.
\end{theorem}
\begin{proof}\leanok
    The standard coordinate equivalences intertwine both actions of $A$,
    and are surjective, giving equality of the ranges. An orthogonal
    projection is the orthogonal projection onto its range. Transport this
    equality through the same standard matrix representation.
\end{proof}

\begin{theorem}[The existing support of vector columns]
    \label{thm:finite_span_support_projection}
    \lean{Matrix.projectionMatrix_span_range_eq_familySupportProj}
    \leanok
    Let $(u_i)$ be a finite family in $\C^m$, and regard each $u_i$ as a
    one-column matrix $C_i:\C\to\C^m$. The existing joint column support
    of $(C_i)$ is exactly the orthogonal projection onto
    $\operatorname{span}\{u_i\}$. No independence or nonzero-vector
    condition is required.
\end{theorem}
\begin{proof}\leanok
    \uses{thm:projection_matrix_range,lem:family_column_support_range}
    The range of the joint support is the range of the matrix whose
    columns are all the $u_i$, hence is their actual span. The joint support
    is an orthogonal projection, so it is the orthogonal projection onto
    this span. The coordinate identification gives the stated equality
    of ambient matrices.
\end{proof}

\begin{theorem}[Borel dependence of the actual span projection]
    \label{thm:borel_finite_span_projection}
    \lean{Matrix.measurable_projectionMatrix_span_range}
    \leanok
    For every finite index set, the map taking a family $(u_i)$ to the
    matrix of the orthogonal projection onto its actual span is Borel
    measurable. This includes rank changes, empty families and zero vectors.
\end{theorem}
\begin{proof}\leanok
    \uses{thm:finite_span_support_projection,lem:borel_family_column_support}
    The one-column matrices depend continuously on the vectors. Apply
    the Borel joint-column-support map and the exact identification of
    Theorem~\ref{thm:finite_span_support_projection}.
\end{proof}


\begin{theorem}[Borel projection onto an arbitrary matrix range]
    \label{thm:borel_arbitrary_matrix_range_projection}
    \lean{Matrix.projectionMatrix_range_eq_familySupportProj_singleton,
        Matrix.measurable_projectionMatrix_range}
    \leanok
    The orthogonal projection onto the range of a rectangular matrix is
    the joint column support of the singleton family consisting of that
    matrix. Consequently this projection is Borel measurable in the matrix,
    including at rank changes and at the zero matrix.
\end{theorem}
\begin{proof}\leanok
    \uses{thm:projection_matrix_range,lem:family_column_support_range,
        lem:borel_family_column_support}
    Both orthogonal projections have the same range. The singleton family
    depends continuously on its matrix, so its existing joint-support
    projection is Borel measurable.
\end{proof}

\begin{theorem}[The true viability event of a matrix range]
    \label{thm:borel_viability_matrix_range}
    \lean{Submodule.measurableSet_isCloseTo_prodExt_map_range}
    \leanok
    Let $P(x)$ be a Borel family of square matrices, let $I$ be a fixed
    unitary coordinate identification, and fix a target subspace $T$ of
    the tensor product with a finite exterior space. For every real
    $\delta$, the event that $I(\operatorname{ran}P(x))$ is
    $\delta$-viable for $T$ is Borel measurable. No projection, rank or
    nonzero assumption is imposed on $P(x)$.
\end{theorem}
\begin{proof}\leanok
    \uses{thm:borel_arbitrary_matrix_range_projection}
    Transport the actual range projection by $I$ and tensor it with the
    exterior identity. Viability is the closed positive-semidefinite
    inequality for the compression to the fixed target projection.
\end{proof}
